
Empirical hypothesis testing is foundational to scientific progress, enabling researchers
to assess how well a sample represents an underlying truth.
A central challenge in this process is determining when enough data have been collected
to make a reliable decision, especially as the accuracy of statistical assessments
depends heavily on sample size.  Careful planning for statistical power and precision
is essential to avoid biased inference and ensure meaningful results \citep{cohen1988}.
% \citet{cohen1988}: ``The power of a statistical test is the probability that it will
% yield statistically significant results.'' Cohen's book established power analysis as
% the standard planning tool for sample-size determination, a problem central to this paper.

Sequential hypothesis testing offers a flexible alternative to fixed-sample designs
\citep{wald1947} by allowing data to be evaluated as it accumulates.
% \citet{wald1947} introduced the Sequential Probability Ratio Test (SPRT), the
% foundational framework for sequential analysis. Wald established that both Type~I and
% Type~II error rates can be controlled while terminating as soon as sufficient evidence
% accumulates --- the principle all subsequent sequential methods, including those studied
% here, build upon.
This approach is widely used in fields such as clinical trials
\citep{jennison2000, fda2010bayesian}, quality control, A/B testing \citep{johari2022}, and financial
monitoring,
% \citet{jennison2000}: Jennison \& Turnbull's textbook establishes group sequential
% methods for clinical trials, where interim analyses at pre-planned time points allow
% early stopping for efficacy or futility while controlling Type~I error via
% alpha-spending functions. This shows that principled sequential stopping has long
% been standard in regulated medical research --- and that naive peeking is explicitly
% prohibited in that literature.
% \citet{fda2010bayesian}: The FDA's 2010 guidance explicitly endorses stopping based
% on ``a sufficiently narrow credible interval'' (\S4.7) and adaptive sequential designs
% with pre-specified interim analyses (\S4.10), reflecting regulatory acceptance of
% precision-based and sequential Bayesian stopping in clinical trials. The guidance
% requires type~I error rate control and pre-specified stopping rules --- requirements
% that DPitG satisfies analytically and without case-by-case simulation.
% \citet{johari2022}: ``A/B tests are typically analyzed via frequentist p-values and
% confidence intervals, but these inferences are wholly unreliable if users endogenously
% choose sample sizes by continuously monitoring their tests.'' Johari et al.\ propose
% always-valid inference using mixture sequential probability ratio tests, ensuring
% valid inference at any stopping time --- a frequentist counterpart to the decoupled
% stopping philosophy advocated in this paper.
offering potential benefits in efficiency, timeliness, and ethical considerations.
Central to this approach is the distinction between two critical components:
the \textit{stopping rule}, which dictates when data collection ends, and the \textit{decision rule}, which determines the final verdict on the hypothesis (e.g., accept, reject, or inconclusive). While these components are conceptually distinct, many operational methods conflate them, triggering a stop exactly when a decision criterion is met.

This simultaneous application of stopping and decision rules is common in widely used
frequentist and Bayesian methods, such as p-value thresholds \citep{fisher1925} or posterior-based heuristics
(e.g., HDI+ROPE \citealp{kruschke2013, kruschke2015doing}, Bayes Factors \citealp{kassraftery1995}).
However, if not carefully designed, this coupling can lead to confirmation bias ---
premature stopping on extreme or unrepresentative early samples --- a phenomenon known
as early peeking \citep{simmons2011}.
% \citet{simmons2011}: ``a researcher who starts with 10 observations per condition
% and then tests for significance after every new per-condition observation finds a
% significant effect 22\% of the time'' even when none exists at the nominal 5\% level.
% This paper quantified how flexible stopping rules inflate false-positive rates,
% directly motivating the decoupled stopping and decision rules discussed here.

Following the \textit{Accuracy in Parameter Estimation} paradigm \citep{maxwell2008},
\cite{kruschke2015doing} advocated for a method that strictly decouples these two steps,
thereby addressing the risk of premature and biased conclusions from early peeking.
He proposed using a predetermined target for posterior precision as the
sole \textit{stopping rule}, independent of the hypothesis outcome. Only once this precision goal
is met is the \textit{decision rule} applied to accept or reject the null hypothesis
(via a posterior credible interval criterion).
% \citet{maxwell2008}:
% Maxwell, Kelley \& Rausch (2008) "Sample Size Planning for Statistical Power and Accuracy
% in Parameter Estimation", Annual Review of Psychology.  This is the foundational paper for
% the AIPE (Accuracy in Parameter Estimation) paradigm, which argues sample sizes should be
% planned to achieve sufficiently narrow confidence/credible intervals --- the direct academic
% precursor to Kruschke's PitG.  Citation IS independently justified; it is not merely
% echoing Kruschke's bibliography.
\cite{kruschke2015doing} demonstrated that this ``Precision is the Goal'' (PitG) method
eliminates the early-peeking bias in outcome sampling. However, our analysis reveals an important
limitation: when the null hypothesis is true, or practically so, PitG often yields a high rate of inconclusive
results, rather than correctly accepting the null.\footnote{\citealt{kruschke2015doing} noted in passing that PitG ``did remain undecided
in some cases'' (\S13.3.2, page 391); our analysis shows this is a systematic and
striking effect, with the sampling budget $N_{\rm max}$ and precision goal $\omega_{\rm goal}$
as the key governing parameters. See below for quantitative examples.}

In this paper, we propose an improved variant we call
``Decisive Precision is the Goal'' (DPitG), in which precision is a necessary but no
longer sufficient stopping condition: data collection continues until \textit{both} the
precision goal and a decisive hypothesis outcome are satisfied simultaneously.
The benefit is most pronounced precisely where PitG's inconclusive rate is highest ---
when the null is true or practically so --- making DPitG particularly valuable for
confirming the null hypothesis.
While scientific incentives often favour the discovery of new effects,
rigorous confirmation of the null is equally critical in many fields: from bioequivalence
trials for generic drugs \citep{karalis2012} to do-no-harm regression testing in software
deployment and futility stopping in programme evaluations.
Unlike \textit{Null Hypothesis Significance Testing} (NHST, see Appendix~\ref{app:nhst}),
which can only fail to reject the null, the precision-based framework enables positive
acceptance as a genuine and actionable verdict --- and DPitG ensures this verdict is
reliably reached.
This fully pre-specified rule is also naturally compatible with pre-registration standards
\citep{nosek2018} and directly addresses the design-stage imprecision implicated in
the replication crisis \citep{osc2015, gelman2014}.

Applied to dichotomous data, this more conservative rule substantially reduces
inconclusive outcomes. In the fair coin case ($\theta_{\rm true}=\theta_{\rm null}=0.5$,
precision goal $\omega_{\rm goal}=0.08$), DPitG reduces PitG's inconclusive rate from 63.3\% to just
2.3\% at a median sampling premium of only 4\% over the precision-goal sample size $N_{\rm goal}$, while preserving a
zero false-positive rate. DPitG's conclusiveness is also robust to the choice of
 $\omega_{\rm goal}$ and holds broadly across true and null parameter values
(Sections~\ref{sec:fair_coin}--\ref{sec:trends}).

Our results are demonstrated through simulations with single-group dichotomous data, though the
framework extends naturally to continuous outcomes (Appendix~\ref{app:expected_stop}) and
two-group comparisons (Appendix~\ref{app:extensions}).
Alongside the empirical results we provide a closed-form planning formula, an online
interactive calculator, and open-source code to support adoption in applied settings.